\documentclass[10pt,a4paper]{amsart}
\usepackage[margin=19mm]{geometry}
\usepackage{amsmath,amssymb,mathtools}
\usepackage[scr=rsfs,cal=euler]{mathalfa}
\usepackage{xfrac}
\usepackage[unicode,hidelinks,bookmarks=false]{hyperref}
\newcommand{\ie}{i.e.\ }
\newcommand{\cf}{cf.\ }
\newcommand{\resp}{respectively\ }
\newcommand{\Subsection}{\subsection}
\providecommand{\nobreakdash}{\nobreak\hbox{-}}
\setlength{\emergencystretch}{2em}
\multlinegap0pt
\usepackage{amssymb}
\usepackage{xcolor}
\newtheorem{lemma}{Lemma}
\newtheorem{remark}{Remark}
\title{Constructing orientable and negative orientable sequences with asymptotically optimal period}
\author{Chris J Mitchell, Peter R Wild}
\date{Source: arXiv:2603.18646v2 (CC BY 4.0)}
\begin{document}
\maketitle

\noindent\emph{Source and licence.} Constructing orientable and negative orientable sequences with asymptotically optimal period. Version arXiv:2603.18646v2 (CC BY 4.0), distributed by arXiv under the Creative Commons Attribution 4.0 International licence, http://creativecommons.org/licenses/by/4.0/, as recorded in the arXiv licence metadata for this exact version and shown on the abstract page https://arxiv.org/abs/2603.18646v2. Full licence notice: \texttt{licenses/arxiv-2603.18646-CC-BY-4.0.txt}.\par\smallskip

\noindent\textit{Editorial scope.} Counted unit: the article's lemma lemma:nega\_circuit\_no\_nega\_tuples (source0.tex lines 879--908). The article's complete original proof of it is delivered with it (source0.tex lines 910--971, one \texttt{proof} environment of 62 lines). No mathematics is written, repaired or re-derived: the statement and the proof are the article's own lines, in the article's own notation, and no retained line is substituted or re-flowed.  Retained uncounted as the source-local context the proof needs: lines 218--228; lines 364--378; lines 657--663; lines 695--710.  Nothing was omitted from any retained span, so the package carries no omission record.  No retained line contains a citation, so the wrapper carries no bibliography and no bibliography entry.  No retained line contains a figure environment, an image-inclusion command or drawing code, so no image is delivered and the package declares no asset dependency.  Editorial formatting only, in addition: an AMS \texttt{amsart} preamble that restates the article's own declarations and macros used by the retained lines, namely the environments \texttt{lemma}, \texttt{remark} on separate counters, together with the standard packages those lines need.  The retained environments print in this document as Lemma 1, Remark 1 in order of appearance, whereas the article prints the same items with the numbering of its own full document.  The complete native source of the article as distributed in the arXiv e-print for this exact version is bundled unchanged as \texttt{sources/196710/source0.tex}.
\begin{lemma}  \label{lemma:nega_moved_by_halfn_is_nega}
Suppose $(a_0,a_1,\dots,a_{n-1})$ is a negasymmetric $n$-tuple.  Then the following $n$-tuples are
also negasymmetric:
\[
\begin{split}
(a_{n/2},a_{n/2+1},\dots,a_{n-1},a_0,a_1,\dots,a_{n/2-1}) & \text{~~~~if $n$ is even, and} \\
(a_{(n+1)/2},a_{(n+1)/2+1},\dots,a_{n-1},x,a_0,a_1,\dots,a_{(n-3)/2}) & \text{~~~~if $n$ is odd.}
\end{split}
\]
where $x$ is either 0 or $k/2$.
\end{lemma}

\begin{lemma}  \label{lemma:nega_for_c}
Suppose $n\geq3$, $k\geq3$, and $c\mid n$.  Then:
\begin{itemize}
\item[(i)] Suppose $(a_0,a_1,\dots,a_{n-1})$ is a $k$-ary $n$-tuple with $a_{i}=a_{i+c}$ for
    every $i$ ($0\leq i\leq n-c-1$).  Then $(a_0,a_1,\dots,a_{n-1})$ is a negasymmetric
    $n$-tuple if and only if $(a_0,a_1,\dots,a_{c-1})$ is a negasymmetric $c$-tuple. Moreover,
    in this case if $c$ is odd then $a_{(c-1)/2}=0$ or $k/2$ (the latter only applying if $k$
    is even).
\item[(ii)] Suppose $[a_0,a_1,\dots,a_{n-1}]$ is a necklace in $B_k(n-1)$ of period dividing
    $c$. Then $[a_0,a_1,\dots,a_{n-1}]$ is a negasymmetric necklace in $B_k(n-1)$ if and only
    if $[a_0,a_1,\dots,a_{c-1}]$ is a negasymmetric necklace in $B_k(c-1)$. Moreover,
    $[a_0,a_1,\dots,a_{n-1}]$ is a negasymmetric necklace in $H_k(n-1)$ if and only if
    $[a_0,a_1,\dots,a_{c-1}]$ is a negasymmetric necklace in $H_k(c-1)$.
\end{itemize}
\end{lemma}

\begin{lemma}  \label{lemma:nega_ciruit_means_nega_tuple_odd_k_odd_c}
Suppose $k\geq3$ and $n\geq 3$.  Suppose $[a_0,a_1,\dots,a_{n-1}]$ is a negasymmetric necklace in
$\mathcal{C}_k(n-1)$ of period $c$, where $c$ is odd ($c\mid n$).  Then
$(a_{\overline{s}},a_{\overline{s+1}},\dots,a_{\overline{s-1}})$ is a negasymmetric $n$-tuple for
some $s$ ($0\leq s\leq n-1$).  Moreover, $[a_0,a_1,\dots,a_{n-1}]$ contains precisely one
negasymmetric $n$-tuple.
\end{lemma}

\begin{remark}  \label{remark:n_even_counterexample}  \label{remark:n_odd_OK}
Since $c\mid n$, Lemma~\ref{lemma:nega_ciruit_means_nega_tuple_odd_k_odd_c} immediately implies
that if $n$ is odd then all the negasymmetric necklaces in $\mathcal{C}_k(n-1)$ contain a unique
negasymmetric $n$-tuple. However, the above result does not hold if $c$ is even, and as a result we
examine the $n$ odd and even cases separately.

To see why the $c$ even case is different, suppose
\[ (a_0,a_1,\dots,a_{n-1}) = (-a_{\overline{t}},-a_{\overline{t-1}},\dots,-a_{\overline{t+1}}) \]
for some $t$.  If $t$ is odd then the same argument as in the above proof shows that there is a
negasymmetric $n$-tuple in the necklace $[a_0,a_1,\dots,a_{n-1}]$.  However, if $t$ is even then
the argument used above does not work, since $c$ being odd is key.  An example in the case $n=4$
and $k=3$ involves the necklace $[a_0=1,a_1=0,a_2=2,a_3=0]$.  It is simple to see that:
\[ (a_0,a_1,a_2,a_3)=(-a_2,-a_1,-a_0,-a_3) \]
i.e.\ the necklace is negasymmetric (in the above terminology $t=2$).  However, none of the four
4-tuples in the necklace are negasymmetric.
\end{remark}

\begin{lemma}  \label{lemma:nega_circuit_no_nega_tuples}
Suppose $k\geq3$, $n\geq 3$, and $[a_0,a_1,\dots,a_{n-1}]$ is a negasymmetric necklace in
$\mathcal{C}_k(n-1)$ containing no negasymmetric $n$-tuples.  Then $n$ is even,
$[a_0,a_1,\dots,a_{n-1}]$ has  period $c$ for some even $c\mid n$, and, for some $s$ ($0\leq s<
c/2$):
\begin{itemize}
\item[(i)] $a_{s}=0$ or $k/2$ and $a_{s+c/2}=0$ or $k/2$;
\item[(ii)] both
\[ (a_{s+1},a_{s+2},\dots,a_{c-1},a_0,a_1,\dots,a_{s-1}) \]
and
\[ (a_{s+c/2+1},a_{s+c/2+2},\dots,a_{c-1},a_0,a_1,\dots,a_{s+c/2-1}) \]
are negasymmetric $(c-1)$-tuples;
\item[(iii)] $[a_0,a_1,\dots,a_{c-1}]$ contains no other negasymmetric
    $(c-1)$-tuples;~moreover, the
    $(n-1)$-tuples of (ii) are distinct;
\item[(iv)]
\[ (a_{s+1},a_{s+2},\dots,a_{n-1},a_0,a_1,\dots,a_{s-1}) \]
and
\[ (a_{s+c/2+1},a_{s+c/2+2},\dots,a_{n-1},a_0,a_1,\dots,a_{s+c/2-1}) \]
are distinct negasymmetric $(n-1)$-tuples,~and
 $[a_0,a_1,\dots,a_{n-1}]$ contains no other negasymmetric $(n-1)$-tuples;~also, if~
 $a_s=a_{s+c/2}$
\[ (a_{s+n/2+1},a_{s+n/2+2},\dots,a_{n-1},a_0,a_1,\dots,a_{s+n/2-1}) \]
 equals either
 \[ (a_{s+1},a_{s+2},\dots,a_{n-1},a_0,a_1,\dots,a_{s-1}) \]
or
\[ (a_{s+c/2+1},a_{s+c/2+2},\dots,a_{n-1},a_0,a_1,\dots,a_{s+c/2-1}) \]
according as $n/c$ is even or odd.
\end{itemize}
\end{lemma}

\begin{proof}
The fact $c$ is even follows immediately from
    Lemma~\ref{lemma:nega_ciruit_means_nega_tuple_odd_k_odd_c}, and since $c\mid n$ it follows that
    $n$ is also even. By Lemma~\ref{lemma:nega_for_c}, $[a_0,a_1,\dots,a_{c-1}]$ is a
    negasymmetric necklace in $B_k(c-1)$ which contains no negasymmetric $c$-tuples. Since
    $[a_0,a_1,\dots,a_{c-1}]$ is negasymmetric, by definition $a_{i}=-a_{\overline{t-i}}$ for
    some $t$ ($0\leq t\leq c-1$) and every $i$ ($0\leq i\leq c-1$), where the bar indicates
    computation modulo $c$.

By Remark~\ref{remark:n_even_counterexample}, we know $t$ is even, or else there would exist a
    negasymmetric $c$-tuple in the necklace. Hence $a_{t/2}=-a_{t/2}$, i.e.\ $a_{t/2}=0$ or $k/2$.  If we
    set $s=t/2\bmod c$, then $a_{s}=a_{t/2}=0$ or $k/2$, since $[a_0,a_1,\dots,a_{n-1}]$ has period $c$.
    Moreover $0\leq s<c/2$ by definition of $s$, establishing part of (i).

It also follows immediately that $a_{s+i}=-a_{\overline{s-i}}$ for every $i$ ($0\leq i\leq
    c-1$), where $\overline{x}=x\bmod c$.  Thus
    $(a_{s+1},a_{s+2},\dots,a_{c-1},a_0,a_1,\dots,a_{s-1})$ is a negasymmetric $(c-1)$-tuple,
    establishing the first part of (ii).

But by Lemma~\ref{lemma:nega_for_c}(i), this means that $a_{s+c/2}=0$ or $k/2$, completing the
proof of (i). By Lemma~\ref{lemma:nega_moved_by_halfn_is_nega}, it follows that
\[ (a_{s+c/2+1},a_{s+c/2+2},\dots,a_{c-1},a_0,a_1,\dots,a_{s+c/2-1}) \]
is negasymmetric and (ii) follows.

Suppose the first part of (iii) is not true, and without loss of generality relabelling subscripts
as necessary, suppose
\[ (a_0,a_1,\dots,a_{c-2}) \text{~~and~~}
(a_{\overline{t}},a_{\overline{t+1}},\dots,a_{\overline{t-2}}) \]
are both negasymmetric
$(c-1)$-tuples where $0<t<c/2$ and $\overline{x}=x\bmod c$.  Thus
\[ a_i=-a_{c-2-i} \text{~~and~~} a_{\overline{t+i}}=-a_{\overline{t-2-i}} \]
for every $i$ ($0\leq i<c-1$).

Also, if $(a_0,a_1,\dots,a_{c-2})$ is negasymmetric then it must have pseudoweight
 $(c-1)k/2$, and thus $a_{c-1}$ has pseudoweight
$k/2$, and hence $a_{c-1}=0$ or $k/2$. Similarly we have $a_{t-1}=0$ or $k/2$. Hence the equations
above also hold for $i=c-1$.

Setting $i=t+j$ in the first equation gives $a_{\overline{t+j}}=-a_{\overline{-2-t-j}}$ for every
$j$ ($0\leq j<c$), and thus $a_{\overline{-t-2-i}}=a_{\overline{t-2-i}}$ for every $i$ ($0\leq
i<c$). Hence $[a_0,a_1,\dots,a_{c-1}]$ has period dividing $2t<c$, which is a contradiction.

The second part of (iii) follows because if the two $(c-1)$-tuples are equal then
$a_{\overline{s+i}}=a_{\overline{s+c/2+i}}$ for every $i$ ($1\leq i\leq c-1$), where the bar
denotes reduction modulo $c$. Hence, setting $i=c/2$, we have
$a_{\overline{s+c/2}}=a_{\overline{s+c}}=a_{\overline{s}}$, and so
$a_{\overline{s+i}}=a_{\overline{s+c/2+i}}$ for every $i$ ($0\leq i\leq c-1$). Thus
$[a_0,a_1,\dots,a_{n-1}]$ has period dividing $c/2$, giving the desired contradiction.



Now consider (iv). That these $(n-1)$-tuples are negasymmetric follows immediately, because
$(a_0,a_1,\dots,a_{n-1})$ is a concatenation of $n/c$ copies of $(a_0,a_1,\dots,a_{c-1})$.
Moreover, another negasymmetric $(n-1)$-tuple contained in $[a_0,a_1,\dots,a_{n-1}]$ would imply
another negasymmetric $(c-1)$-tuple contained in $[a_0,a_1,\dots,a_{c-1}]$. That they
are distinct follows from exactly the same argument used to show the $(c-1)$-tuples of (ii) are
distinct, except this time the subscripts are computed modulo $n$.

The statement regarding
\[ (a_{s+n/2+1},a_{s+n/2+2},\dots,a_{n-1},a_0,a_1,\dots,a_{s+n/2-1}) \]
is immediate by the periodicity of $[a_0,a_1,\dots,a_{n-1}]$.
\end{proof}

\end{document}
